\documentclass[10pt,a4paper]{article}
\usepackage[T1]{fontenc}
\usepackage[margin=23mm]{geometry}
\usepackage{amsmath,amssymb,amsthm,mathtools,xcolor,booktabs,array}
\usepackage{tikz-cd,xurl}
\usepackage[hidelinks]{hyperref}
\hypersetup{pdftitle={A representation-theoretic proof of the AMS relation},
 pdfauthor={John Martin},
 pdfsubject={Lean-certified AMS coefficient proof and one shared actual vertex table for both families}}
\newtheorem{theorem}{Theorem}
\newtheorem{lemma}[theorem]{Lemma}
\newtheorem{corollary}[theorem]{Corollary}
\newcommand{\Sym}{\operatorname{Sym}}
\newcommand{\Inv}{\operatorname{Inv}}
\newcommand{\Hom}{\operatorname{Hom}}
\newcommand{\tr}{\operatorname{tr}}
\newcommand{\qf}[1]{[#1]_v!}
\newcommand{\qd}[1]{[#1]_v}
\newcommand{\id}{\mathbf 1}
\newcommand{\lean}[1]{{\small\nolinkurl{#1}}}
\newenvironment{revision}{\begingroup}{\endgroup}
\setlength{\emergencystretch}{1em}
\newcommand{\brac}[2]{\begin{Bmatrix}#1\\#2\end{Bmatrix}}
% aiXiv local preparation copy. No submission has been made.
% Author and declaration metadata confirmed by the user on 25 November 2026.
\title{A representation-theoretic proof of the AMS relation}
\author{John Martin}
\date{25 November 2026\\[0.5em]
  {\small AI-derived proof and AI-written manuscript under the author's direction.}}
\begin{document}
\begin{revision}
\maketitle
\begin{abstract}
We give a representation-theoretic proof of the factorial formula for
multiplicity-free quantum $6j$ symbols in the two symmetric families
stated in Eq.~(35) of V.~Alekseev, A.~Morozov and A.~Sleptsov,
\emph{Multiplicity-free $U_q(\mathfrak{sl}_N)$ $6$-$j$ symbols:
Relations, asymptotics, symmetries}, \emph{Nuclear Physics B}
\textbf{960}, 115164 (2020),
\href{https://doi.org/10.1016/j.nuclphysb.2020.115164}{doi:10.1016/j.nuclphysb.2020.115164}.
For real $v>0$, $v\ne1$, the proof constructs the physical coupling
spaces, their compact forms, the complete irreducible channels, and
an explicit comparison with rank-two coupling spaces. In Type II,
a universal normalization and a full shifted-channel comparison fix
the exact coefficient formula. For $N\ge3$, we construct one table
of full isometric intertwiners indexed by actual ordered Dynkin-weight
triples and shared by both families. Its Type-I trees agree with the
calibrated trees; both Type-II trees acquire the same exterior sign,
which cancels in their overlap. The construction includes coincident
mixed and two-row carriers at $N=3$ and all zero-degree endpoints.
We also establish the physical coefficient deformation limits at
$v=1$ and equality of compatible complex analytic continuations.
The shared vertex table is constructed at real $v>0$, $v\ne1$.
The finite proof and the coefficient continuation results are certified
in Lean 4.24.0 with Mathlib v4.24.0.
\end{abstract}
\end{revision}

\begin{revision}
\noindent\textbf{Certified scope and convention.}
Nawata, Ramadevi and Zodinmawia proposed a factorial formula for
multiplicity-free quantum $6j$ symbols with symmetric labels~\cite{NRZ}.
Alekseev, Morozov and Sleptsov (AMS) rewrote it as Eq.~(35) and related it
to the rank-two factorial evaluation in Eq.~(50)~\cite{AMS}. Here we prove
that the formula evaluates the actual coupling coefficients of both families
for all $N\ge2$ and all admissible exterior and channel labels.
The direct finite proof appears in Section~\ref{sec:certified}. It
constructs the higher-rank channel vectors independently of the factorial
formula and the Kazhdan--Lusztig equivalence. Section~\ref{sec:KZroute}
gives a second proof of rank reduction via the classical determinant
model and the KZ associator. Section~\ref{sec:phases} fixes the signs,
including a common vertex table at $N\ge3$, and
Section~\ref{sec:continuation} states the classical coefficient limit and
complex continuation. The Lean statements and replay are recorded in
Section~\ref{sec:audit}.
\end{revision}

\medskip
Set $h=N-2$, $N\ge2$, and write $S_r^{(N)}=\Sym^r\mathbb C^N$.
For $v>0$, $v\ne1$, our quantum integer and factorial are
\[
 [k]_v=\frac{v^k-v^{-k}}{v-v^{-1}},\qquad
 \qf{k}=\prod_{a=1}^k[a]_v\quad(k\ge0).
\]
At $v=1$ we use their continuous values $[k]_1=k$ and $[k]_1!=k!$.
Thus the parameter called $q$ in AMS, with half-powers in its quantum integers,
is $v^2$. The parameter of $U_v(\mathfrak{sl}_N)$ below is $v$.

\section{Statement for the actual coupling matrices}
Consider the four exterior representations
\begin{align}
\mathrm I:\quad&
 S_{r_1}^{(N)}\otimes(S_{r_2}^{(N)})^*\otimes
 S_{r_3}^{(N)}\otimes(S_{r_4}^{(N)})^*,
 &r_1+r_3&=r_2+r_4,\label{eq:typeI}\\
\mathrm {II}:\quad&
 S_{r_1}^{(N)}\otimes S_{r_2}^{(N)}\otimes
 (S_{r_3}^{(N)})^*\otimes(S_{r_4}^{(N)})^*,
 &r_1+r_2&=r_3+r_4.\label{eq:typeII}
\end{align}
At $v\ne1$ use the corresponding type-1 quantum modules. In each
family let $V_1,V_2,V_3,V_4$ denote the four factors in the displayed order.
Identify the singlet space with $\Hom(V_4^*,V_1\otimes V_2\otimes V_3)$
by duality. Write $L_i$ and $R_j$ for normalized coupling vectors, viewed
in the same space by the tensor associator. Our overlap matrix has entries
$\mathsf U_{N,T;ij}(v)=\langle L_i,R_j\rangle$, with rows indexed by the
$12$ channel and columns by the $23$ channel. The inner product is
conjugate-linear in its first argument. Thus its adjoint acts on coordinate
columns from the $12$ basis to the $23$ basis.

The channel labels are
\begin{align*}
 i&=\max(|r_1-r_2|,|r_3-r_4|),\ldots,
       \min(r_1+r_2,r_3+r_4),\\
 j&=\max(|r_1-r_4|,|r_2-r_3|),\ldots,
       \min(r_1+r_4,r_2+r_3),
\end{align*}
in steps of two. For a mixed pair with symmetric degree $a$ and dual
symmetric degree $b$, contraction $u$ gives highest weight
$(a-u)\omega_1+(b-u)\omega_{N-1}$ and numerical label $a+b-2u$.
For a symmetric pair the label $a+b-2t$ specifies the Young diagram
$(a+b-t,t)$. A pair of dual symmetric modules gives the dual of that
representation. At $N=2$ the numerical label in either case is the
usual highest weight. The Pieri and mixed harmonic decompositions are
multiplicity-free. Matching dual channels gives a one-dimensional space
for each admissible channel in the singlet space.

\begin{theorem}[Physical rank reduction]\label{thm:physical}
For $v>0$ and every admissible set of labels in \eqref{eq:typeI} or
\eqref{eq:typeII},
\begin{equation}\label{eq:physical}
 \mathsf U_{N,T}(v)=D_L(v)\,
 \mathsf U_{2}(v;\boldsymbol r';i',j')\,D_R(v),
\end{equation}
where $D_L,D_R$ are diagonal unitary phase matrices and
\begin{center}
\begin{tabular}{c@{\qquad}cccccc}
\toprule
$T$ & $r'_1$ & $r'_2$ & $r'_3$ & $r'_4$ & $i'$ & $j'$\\
\midrule
I & $r_1$ & $r_2+h$ & $r_3$ & $r_4+h$ & $i+h$ & $j+h$\\
II & $r_1+h$ & $r_2+h$ & $r_3$ & $r_4$ & $i$ & $j+h$\\
\bottomrule
\end{tabular}
\end{center}
\begin{revision}
For $v>0$, $v\ne1$, the direct quantum proof below establishes the theorem
on complete physical highest spaces. The coefficient limit at $v=1$ and
all compatible complex analytic continuations are certified separately
as stated in Section~\ref{sec:continuation}. The retained classical argument
also establishes the classical representation-theoretic rank reduction.
Unitarity concerns positive real parameters. For arbitrary real normalized
coupling bases, the comparison factors at each parameter are whole-row
and whole-column signs. Exact phase choices are specified in
Section~\ref{sec:phases}.
\end{revision}
\end{theorem}

\section{One polynomial model for both families}
Use vectors $x,z$ and covectors $y,w$, and write
\[
 Z=\begin{pmatrix}x\cdot y&x\cdot w\\z\cdot y&z\cdot w\end{pmatrix}.
\]
Let $(A,B)$ be its row degrees and $(C,D)$ its column degrees:
\[
\begin{array}{c|cccc}
T&A&B&C&D\\\hline
\mathrm I&r_1&r_3&r_2&r_4\\
\mathrm {II}&r_1&r_2&r_3&r_4
\end{array}
\qquad A+B=C+D=\rho.
\]
Write $\mathcal H_{N,T}$ for the corresponding invariant space.

\subsection{A basis from a dimension count}
\begin{lemma}\label{lem:pairingbasis}
Put $b=\min(A,B,C,D)$. The space $\mathcal H_{N,T}$ has dimension $b+1$
and the following basis:
\begin{equation}\label{eq:monomial}
 e_k=Z_{11}^{k}Z_{12}^{A-k}Z_{21}^{C-k}Z_{22}^{B-C+k},
 \qquad \ell=\max(0,C-B)\le k\le u=\min(A,C).
\end{equation}
\end{lemma}
\begin{proof}
Pieri's multiplicity-free decomposition \cite[Sec.~6.1, Eq.~(6.8)]{FH} is
\begin{equation}\label{eq:pieri}
 S_A^{(N)}\otimes S_B^{(N)}
 =\bigoplus_{t=0}^{\min(A,B)}\mathbb S_{(\rho-t,t)}\mathbb C^N.
\end{equation}
The analogous $C,D$ decomposition and Schur's lemma give exactly the
common channels $0\le t\le b$, hence dimension $b+1$.
Among diagrams of the same total degree, equality of the $SL_N$ highest
weights implies equality of the diagrams. Their row lengths could differ
only by a constant vector, and the common total degree makes it zero. Each $e_k$
is invariant. They are independent because setting $x=e_1,z=e_2$ makes
all four entries of $Z$ independent variables, in every $N\ge2$.
Finally $u-\ell+1=b+1$, so these vectors form a basis.
\end{proof}
This uses the same pair decomposition already needed for channel
multiplicity-freeness, and avoids a separate invariant-algebra theorem.

\subsection{Determinant transport, including surjectivity}
In dimension two put $D_x=\det(x,z)$ and $D_y=\det(y,w)$ and define
\begin{equation}\label{eq:Fmap}
 F_{\mathrm I}f=D_y^h f(Z),\qquad F_{\mathrm {II}}f=D_x^h f(Z).
\end{equation}
Write $\mathcal H_2(\boldsymbol r')$ for the shifted $SL_2$ invariant
space. These maps have exactly the exterior degrees in
Theorem~\ref{thm:physical}.
\begin{lemma}\label{lem:bijection}
Both maps are bijections onto the shifted $SL_2$ invariant spaces.
\end{lemma}
\begin{proof}
Independence of $Z$ and multiplication by a nonzero determinant prove
injectivity. For Type I, the row pair has $SL_2$ labels $\rho-2t$ and
the shifted column pair has $\rho+2h-2s$. Matching gives $s=t+h$,
with precisely $0\le t\le b$. Type II interchanges which pair is
shifted and gives the same count. Thus the target also has dimension
$b+1$, proving surjectivity.
\end{proof}

\subsection{Exact isometry by creation--annihilation adjunction}
Write $(X_1,X_2)=(x,z)$ and $(Y_1,Y_2)=(y,w)$ for the two vector and
two covector \emph{slots}. Their components still range from $1$ to $N$.
Let $M_{ij}$ multiply by $Z_{ij}=X_i\cdot Y_j$ and let
$\Delta_{ij}^{(N)}=\sum_\alpha
\partial_{(X_i)_\alpha}\partial_{(Y_j)_\alpha}$.
For a polynomial $f$ in the four formal entries of $Z$, the chain rule is
\begin{equation}\label{eq:radial}
 \Delta_{ij}^{(N)}f(Z)=
 N\partial_{ij}f+\sum_{a,b=1}^{2}Z_{ab}\partial_{ib}\partial_{aj}f.
\end{equation}
For $J=D_x$ or $D_y$, the two-dimensional Euler identities for a
determinant give, for all four pairs $i,j$,
\begin{equation}\label{eq:detderivative}
 \Delta_{ij}^{(2)}(J^hf(Z))
 =J^h\bigl(\Delta_{ij}^{(2)}f(Z)+h\partial_{ij}f(Z)\bigr).
\end{equation}
For example $x\cdot\partial_xD_x=z\cdot\partial_zD_x=D_x$ and
$z\cdot\partial_xD_x=x\cdot\partial_zD_x=0$; the covector identities
are identical. Since $2+h=N$, we have the stronger transport identities
\begin{equation}\label{eq:laddertransport}
 M_{ij}^{(2)}F_T=F_TM_{ij}^{(N)},\qquad
 \Delta_{ij}^{(2)}F_T=F_T\Delta_{ij}^{(N)}.
\end{equation}
Use the same definition of $F_T$ in each balanced multidegree: these
identities relate neighboring homogeneous spaces.

\begin{lemma}[Universal normalization]\label{lem:isometry}
For the standard Fock inner product
$\langle X^\alpha,X^\beta\rangle=\delta_{\alpha\beta}\alpha!$,
\begin{equation}\label{eq:scalarisometry}
 \langle F_Tf,F_Tg\rangle_2=h!(h+1)!\,\langle f,g\rangle_N,
 \qquad
 \boxed{S_T=\frac{F_T}{\sqrt{h!(h+1)!}}\ \text{is unitary}.}
\end{equation}
The normalization is independent of the family and all exterior labels.
\end{lemma}
\begin{proof}
Put $c_h=h!(h+1)!$. In pairing degree zero, $F_T1=J^h$.
The $h+1$ monomials in the binomial expansion of $J^h$ are orthogonal;
each contributes $(h!)^2$. Hence $\|J^h\|^2=(h+1)(h!)^2=c_h$.
Induct on the common pairing degree of $f,g$. By
Lemma~\ref{lem:pairingbasis} it suffices to take $f$ a pairing monomial,
so at positive degree $f=M_{ij}f_0$ for some $i,j$.
Fock adjunction, \eqref{eq:laddertransport} and induction give
\begin{align*}
 \langle F_T(M_{ij}f_0),F_Tg\rangle_2
 &=\langle F_Tf_0,\Delta_{ij}^{(2)}F_Tg\rangle_2\\
 &=\langle F_Tf_0,F_T\Delta_{ij}^{(N)}g\rangle_2\\
 &=c_h\langle f_0,\Delta_{ij}^{(N)}g\rangle_N
 =c_h\langle M_{ij}f_0,g\rangle_N.
\end{align*}
Here \eqref{eq:radial} expresses $\Delta_{ij}^{(N)}g$ as a polynomial
of pairing degree one less, in the required neighboring multidegree.
Sesquilinearity completes the induction. Lemma~\ref{lem:bijection} makes the
resulting isometry onto.
\end{proof}

\subsection{The adjacent Casimirs and their channel labels}
Set $P_N=M_{11}\Delta_{11}^{(N)}$ and
$Q_N=M_{21}\Delta_{21}^{(N)}$. Equation~\eqref{eq:laddertransport} gives
\begin{equation}\label{eq:mixedtransport}
 P_2S_T=S_TP_N,\qquad Q_2S_T=S_TQ_N.
\end{equation}
Normalize roots to squared length two and write $\Omega$ for the split
Casimir, $\tfrac12(C_{V\otimes W}-C_V\otimes\id-\id\otimes C_W)$
for a pair of modules $V,W$.
For the mixed pair $x,y$ of degrees $a,b$, the $\mathfrak{gl}_N$ split
Casimir is $-P_N$ since the covector action is the negative transpose.
The central term is $-ab/N$, so
\begin{equation}\label{eq:mixedomega}
 \Omega_{xy}^{(N)}=-P_N+\frac{ab}{N}\id.
\end{equation}
For other mixed pairs the same formula uses their own $M_{ij}\Delta_{ij}^{(N)}$.
For the same-orientation pair put
$E=x\cdot\partial_z$, $G=z\cdot\partial_x$ and $E_x=x\cdot\partial_x$.
Its $\mathfrak{gl}_N$ split Casimir is $EG-E_x$. The corresponding
$\mathfrak{sl}_N$ operator subtracts the central term $ab/N$.
Both $E$ and $G$ annihilate $D_x$, and the Leibniz rule gives
$E_x(D_x^h f)=D_x^h(E_x f+hf)$. Hence
\begin{equation}\label{eq:same}
 \Omega_{xz}^{\mathfrak{gl}_2}S_{\mathrm {II}}
 =S_{\mathrm {II}}\Omega_{xz}^{\mathfrak{gl}_N}-hS_{\mathrm {II}}.
\end{equation}
Equations \eqref{eq:mixedomega} and \eqref{eq:same} give
\begin{equation}\label{eq:affine}
 \Omega_{12}^{(2)}S_T=S_T\Omega_{12}^{(N)}+c_{12}S_T,
 \qquad
 \Omega_{23}^{(2)}S_T=S_T\Omega_{23}^{(N)}+c_{23}S_T.
\end{equation}
Explicitly, with $(a,b,c,d)=\boldsymbol r$,
\begin{align*}
 \mathrm I:\quad&c_{12}=\frac{a(b+h)}2-\frac{ab}{N},
 &c_{23}&=\frac{c(b+h)}2-\frac{bc}{N},\\
 \mathrm {II}:\quad&c_{12}=-h-\frac{(a+h)(b+h)}2+\frac{ab}{N},
 &c_{23}&=\frac{(b+h)c}2-\frac{bc}{N}.
\end{align*}
On a mixed channel of contraction number $s$, the eigenvalue of
$M\Delta$ is
\begin{equation}\label{eq:mixedspec}
 s(a+b+N-s-1),\qquad 0\le s\le\min(a,b).
\end{equation}
This follows by iterating $\Delta M-M\Delta=E_x+E_y+N$ on harmonics.
The eigenvalues are distinct. The determinant shift increases one exterior
degree by $h$ and keeps $s$ fixed, so the mixed label increases by $h$.
For the same-orientation channel $(\rho-t,t)$, $\rho=a+b$,
\begin{equation}\label{eq:samespec}
 \Omega_{xz}^{\mathfrak{gl}_N}=ab-t(\rho-t+1).
\end{equation}
Substitution of $(a+h,b+h,t+h)$ decreases this by $h$, exactly as in
\eqref{eq:same}, and leaves the label $a+b-2t$ unchanged.
The lower-bound calculation uses
$\max(|x+h|,|x-h|)=|x|+h$ for $h\ge0$.
The shifted triangle bounds give precisely these channels: Type I
translates both ranges by $h$; Type II keeps the $i$ range and translates
the $j$ range. Thus $S_T$ matches both normalized classical coupling
bases, up to phases, with exactly the labels in the theorem.

\section{An independent analytic route to the quantum statement}\label{sec:KZroute}
\begin{revision}
This section retains the analytic argument based on the published unitary
Kazhdan--Lusztig equivalence. That external theorem is not a dependency or
hypothesis of the completed Lean coefficient certificate: the direct
quantum proof is given in Section~\ref{sec:certified}.
\end{revision}
\begin{lemma}[Affine Casimir transport of recoupling]\label{lem:KZ}
Suppose two classical multiplicity spaces have one-dimensional channel
spaces and a unitary bijection $S$ satisfying
$A'S=S A+\alpha S$ and $B'S=S B+\beta S$ for their two pair Casimirs,
with real $\alpha,\beta$. Suppose also that $S$ matches the specified
channel labels. Their quantum recoupling matrices, at the same positive
deformation parameter, then agree up to row and column phases.
\end{lemma}
\begin{proof}
Put $\eta=\log v/(\pi i)\in i\mathbb R$. The KZ associator is the
regularized connection matrix for
\begin{equation}\label{eq:KZ}
 \frac{dG}{dz}=\eta\left(\frac A z+\frac B{z-1}\right)G.
\end{equation}
Write $G_0$ for the solution normalized at $0$ and $G_1$ for that normalized
at $1$, so $\Phi(\eta A,\eta B)=G_1^{-1}G_0$. Simultaneous conjugation by
$S$ conjugates both solutions. Adding $\alpha\id$ to $A$ and
$\beta\id$ to $B$ multiplies both normalized solutions by the same scalar
$z^{\eta\alpha}(1-z)^{\eta\beta}$. This scalar cancels from their ratio.
Therefore
\begin{equation}\label{eq:Phiintertwine}
 \Phi(\eta A',\eta B')S=S\Phi(\eta A,\eta B).
\end{equation}
These properties are also stated directly in \cite[Sec.~2, p.~14]{NT}.

The nontrivial representation-theoretic input is the \emph{unitary}
Kazhdan--Lusztig/Drinfeld equivalence, in the precise form of
\cite[Theorem~3.1]{NT}: for every $v>0$ the category of classical compact
group modules with associator $\Phi(\eta\Omega_{12},\eta\Omega_{23})$
is unitarily monoidally equivalent to the category of $U_v(\mathfrak g)$
modules, with irreducible highest weights identified. Thus a quantum
recoupling matrix is the matrix of this associator between the two
\emph{classical} orthonormal coupling bases, up to their vertex phases.
In the classical invariant-tensor realization, bending an intertwiner
with the unnormalized coevaluation gives the Hilbert--Schmidt pairing
of the two intertwiners. If $f^*g=\langle f,g\rangle\id$ on the output
irrep, this pairing is $\dim(V_4^*)\langle f,g\rangle$. The Fock metric
can supply an additional fixed scalar for the exterior degrees. Both
scalars cancel on normalization. In the convention of the overlap matrix
above, the associator gives entries
\begin{equation}\label{eq:Fentry}
 \langle \Phi(\eta A,\eta B)L_i,R_j\rangle
\end{equation}
up to row and column phases. The opposite pairing gives the conjugate
entries of the corresponding coordinate change. The tensor equivalence
is unitary, so the choices on the one-dimensional vertex spaces differ
only by unit phases.

Since $S$ sends $L_i$ and $R_j$ to their matched channel morphisms up to
phases, \eqref{eq:Phiintertwine} makes the entries \eqref{eq:Fentry} equal
up to independent row and column phases. This proves the lemma.
\end{proof}

\begin{proof}[Proof of Theorem~\ref{thm:physical}]
Apply Lemma~\ref{lem:KZ} with $A=\Omega_{12}^{(N)}$,
$B=\Omega_{23}^{(N)}$, and $S=S_T$ from \eqref{eq:scalarisometry}.
Lemmas~\ref{lem:bijection} and \ref{lem:isometry}, together with
\eqref{eq:affine}--\eqref{eq:samespec}, verify every hypothesis.
In diagram form, the established commuting square is
\[
\begin{tikzcd}[column sep=large,row sep=large]
\mathcal H_{N,T} \arrow[r,"S_T"]
 \arrow[d,"{\Phi(\eta\Omega^{(N)}_{12},\eta\Omega^{(N)}_{23})}"']
 &\mathcal H_{2}(\boldsymbol r')
 \arrow[d,"{\Phi(\eta\Omega^{(2)}_{12},\eta\Omega^{(2)}_{23})}"]\\
\mathcal H_{N,T} \arrow[r,"S_T"']&\mathcal H_{2}(\boldsymbol r').
\end{tikzcd}
\]
The normalized recoupling matrices therefore agree up to vertex phases
for every $v>0$. At $v=1$ the associator is the identity and the comparison
reduces to the classical isometry $S_T$. Section~\ref{sec:continuation}
constructs the continuation of the calibrated coefficient matrix from
explicit finite expressions.
\end{proof}

\section{Recovery of the complete signed AMS sum}
Let $X_i,Y_j$ be the two higher-rank intermediate irreducibles and let
$d_i=\dim_v X_i$, $e_j=\dim_v Y_j$. Define
\begin{align}
 \theta_N(a,b,c)^2
 &=\frac{\qf{(a+b-c)/2}\qf{(a+c-b)/2}\qf{(b+c-a)/2}}
          {\qf{(a+b+c)/2+N-1}},\label{eq:theta}\\
 K_N&=\qf{N-1}\qf{N-2}
 \theta_N(r_1,r_2,i)\theta_N(r_3,r_4,i)
 \theta_N(r_1,r_4,j)\theta_N(r_2,r_3,j).\label{eq:KN}
\end{align}
At positive $v$ all square roots in this section are positive.
$K_2'$ denotes the same expression at the shifted rank-two labels.

\begin{lemma}[Normalization cancellation]\label{lem:norm}
For either type,
\begin{equation}\label{eq:norm}
 d_i e_j\left(\frac{K_N}{K_2'}\right)^2
       =[i'+1]_v[j'+1]_v.
\end{equation}
\end{lemma}
\begin{proof}
Put $c_N=\qf{N-1}\qf{N-2}$.
The quantum Weyl dimension formula gives, for a mixed channel $(p,q)$,
\begin{equation}\label{eq:mixeddim}
 D_N(p,q)=\frac{[p+q+N-1]_v\qf{p+h}\qf{q+h}}
                    {\qf p\qf q\,c_N},
\end{equation}
and for a two-row channel $(\rho-t,t)$,
\begin{equation}\label{eq:samedim}
 D_N(\rho-t,t)=
 \frac{[\rho-2t+1]_v\qf{\rho-t+h+1}\qf{t+h}}
 {\qf{\rho-t+1}\qf t\,c_N}.
\end{equation}
For Type I, at the $12$ face write $u=(r_1+r_2-i)/2$,
$p=r_1-u$, $q=r_2-u$. The ratio of the squared triangle factors at this
face is $\qf q/\qf{q+h}$, and at the complementary $34$ face it is
$\qf p/\qf{p+h}$. Their product times \eqref{eq:mixeddim} is
$[i+h+1]_v/c_N$. The same calculation on the other two faces gives
$[j+h+1]_v/c_N$.

For Type II put $\rho=r_1+r_2=r_3+r_4$ and
$t=(\rho-i)/2$. The ratios at faces $12,34$ are respectively
\[
 \frac{\qf t}{\qf{t+h}},\qquad
 \frac{\qf{\rho-t+1}}{\qf{\rho-t+h+1}}.
\]
Multiplication by \eqref{eq:samedim} leaves $[i+1]_v/c_N$.
The $23,14$ faces are mixed and leave $[j+h+1]_v/c_N$, as above.
The factor $c_N^2$ from \eqref{eq:KN} cancels both denominators.
\end{proof}

Set $\rho=(r_1+r_2+r_3+r_4)/2$ and
\begin{align*}
 a_0&=(r_3+r_4-i)/2,&b_0&=(r_1+r_2-i)/2+h\delta_{T,\mathrm {II}},\\
 c_0&=(r_2+r_3-j)/2,&d_0&=(r_1+r_4-j)/2,\\
 e_0&=(i+j-r_2-r_4)/2,&f_0&=(i+j-r_1-r_3)/2+h\delta_{T,\mathrm I}.
\end{align*}
Define the fully supported factorial sum
\begin{equation}\label{eq:AMSkernel}
 I_{N,T}=\sum_{u=\max(0,-e_0,-f_0)}^{\min(a_0,b_0,c_0,d_0)}
 \frac{(-1)^{\rho-u}\qf{\rho+N-1-u}}
 {\qf u\qf{a_0-u}\qf{b_0-u}\qf{c_0-u}\qf{d_0-u}
  \qf{e_0+u}\qf{f_0+u}}.
\end{equation}
Every factorial argument in the sum is nonnegative. These offsets give
precisely the support of AMS Eq.~(35). In Type I the balance gives
$f_0=e_0+h$, so the lower bound is $\max(0,-e_0)$. In Type II it gives
$b_0=a_0+h$, so the upper bound is $\min(a_0,c_0,d_0)$. Thus neither
shifted offset adds summation terms relative to the bounds of AMS.

\begin{corollary}[AMS Eq.~(35)]\label{cor:AMS}
For $v>0$, $v\ne1$, choose the channel phases as in
Section~\ref{sec:phases} for each admissible exterior quadruple. The physical higher-rank
quantum $6j$ symbol for either family is
\begin{equation}\label{eq:AMSproved}
 \boxed{\displaystyle\Phi_{N,T}(\boldsymbol r;i,j)=K_N I_{N,T}.}
\end{equation}
Equivalently,
\begin{equation}\label{eq:rank6j}
 \Phi_{N,T}=(-1)^h\sqrt{\frac{[i'+1]_v[j'+1]_v}{d_i e_j}}
 \brac{r'_1/2&r'_2/2&i'/2}{r'_3/2&r'_4/2&j'/2}_{v}.
\end{equation}
For preassigned normalized vertex conventions, these formulas carry
the corresponding independent row and column phase factors.
\end{corollary}
\begin{proof}
First apply the independently established physical reduction
\eqref{eq:physical}. The ordinary rank-two quantum Racah evaluation
\cite{KR} has triangle product $K_2'$ and scalar factorial sum $I_2'$.
\begin{revision}
Both the physical reduction and this finite rank-two evaluation are now
proved in Lean by the direct construction in Section~\ref{sec:certified}.
The certified phase convention is made explicit in Section~\ref{sec:phases}.
\end{revision}
Under the shifts in Theorem~\ref{thm:physical},
$\rho'=\rho+h$ and all six rank-two denominator offsets are exactly
$(a_0,b_0,c_0,d_0,e_0,f_0)$. The numerator factorial is unchanged because
$\rho'+1-u=\rho+N-1-u$. Thus, including the entire effective support,
\begin{equation}\label{eq:scalarshift}
 I_{N,T}=(-1)^h I_2',\qquad
 \brac{r'_1/2&r'_2/2&i'/2}{r'_3/2&r'_4/2&j'/2}_{v}=K_2'I_2'.
\end{equation}

Lemma~\ref{lem:norm} shows that the positive normalization needed to
convert the physical normalized matrices is exactly $K_N/K_2'$.
The only remaining freedom is therefore the phases already present in
Theorem~\ref{thm:physical}. Fix them by transporting the standard rank-two
Racah convention through that theorem and choosing the common sign
$(-1)^h$ for the $6j$ coefficient. For the fixed exterior quadruple this chooses channel
vertex phases, independently of individual matrix entries. It defines the
AMS-compatible convention in \eqref{eq:rank6j}. With the standard rank-two
recoupling convention
$\mathsf U_{2,ij}=(-1)^{\rho'}\sqrt{[i'+1]_v[j'+1]_v}\{6j\}_{v}$,
it amounts to $\sqrt{d_i e_j}\Phi_{N,T}=(-1)^\rho\mathsf U_{2,ij}$.
Finally \eqref{eq:scalarshift} gives \eqref{eq:AMSproved}.

\begin{revision}
The logical order in both routes is the same: first identify the actual
higher-rank overlaps by independent physical transport (the analytic
argument in Section~\ref{sec:KZroute} or the direct certified argument in
Section~\ref{sec:certified}), then evaluate the factorial sum.
The physical reduction is proved without the higher-rank candidate.
\end{revision}
\end{proof}

\begin{revision}
\section{The complete finite quantum proof}\label{sec:certified}
The following is the direct proof used by the Lean certificate. It proves
the physical rank reduction independently of the AMS candidate and of the
Kazhdan--Lusztig equivalence. The classical determinant argument above is
retained as a separate analytic proof and explanation of the shifts.
Throughout this section $v>0$, $v\ne1$, and $N=h+2$.

\subsection{Physical spaces, metrics and completeness}
The source spaces are actual highest-weight multiplicity spaces in
\[
 \mathrm I:\quad S_l\otimes S_m^*\otimes S_c\longleftarrow S_d,
 \quad d=l+c-m\ge0,
 \qquad
 \mathrm {II}:\quad S_a\otimes S_b\otimes S_c^*\longleftarrow S_d,
 \quad a+b=c+d.
\]
They are realized by finite occupation modules. The generator actions are
checked against the defining quantum-group relations; tensor actions,
duality maps, compact positive metrics and the relevant intertwiners are
constructed explicitly. The highest spaces are defined by their raising
and weight equations, before any recoupling formula is introduced.

For Type I put
\begin{equation}\label{eq:certIindices}
 \mathcal J(l,m,c)=\{s\in\mathbb Z:\max(0,m-c)\le s\le\min(l,m)\}.
\end{equation}
The physical cup vectors $b_s$ have exactly these indices. Their selected
occupation coordinates are triangular with nonzero diagonal, proving
linear independence. The actual mixed pair decompositions and Hom bending
give the opposite dimension bound for the whole highest space: the common
channels obey $r-s=c-m$, hence precisely \eqref{eq:certIindices}.
Thus the cup span equals the entire independently defined highest space.
The same argument proves completeness after the rank-two shift.

For Type II let
\begin{equation}\label{eq:certIIindices}
 B=\min(a,b,c,d),\qquad r=0,\ldots,B,
 \qquad t\in\mathcal J(a,c,b).
\end{equation}
The actual two-row channel intertwiners, bent with the physical cup and
divided by their positive norms, give orthonormal vectors $L_r$.
The complete two-row decompositions on the two sides give dimension
$B+1$; therefore these vectors form a basis of the entire highest space.
The independently constructed mixed right channels give a complete
orthonormal basis $R_t$ of the same space. Their physical channel membership
is proved using the actual intertwiners, rather than postulated from an
overlap matrix. These completeness results include all zero labels and
one-dimensional boundary cases.

\subsection{Type I: simultaneous radial transport}
Let $P_N,Q_N$ be the two actual cup creation--annihilation radial operators
on the complete Type-I highest space. Their matrices in the cup basis
are opposite bidiagonal matrices, including the vanishing endpoint terms.
Here is the literal coefficient calculation. Write $n=N-1$ and let
$\beta$ denote the middle background degree, to distinguish it from the
rank shift $h$. The cup monomials $b_{s,k}$ have $s+k=m$, with total
middle exterior degree $\beta+m$. The left operator acts by
\begin{align}
 P_{n,\beta}b_{s,k}
   &=A_{n,\beta}(s,k)b_{s,k}
       +B_{n,\beta}(s,k)b_{s+1,k-1},\notag\\
 A_{n,\beta}(s,k)
   &=v^{-l-\beta-s-k-n+2}[s]_v[l+\beta+k+n]_v,\notag\\
 B_{n,\beta}(s,k)
   &=(-1)^n v^{1-n-\beta-c-2s}[l-s]_v[k]_v.
 \label{eq:certIcoefficients}
\end{align}
The right operator is obtained by interchanging $(l,s)$ with $(c,k)$.
A forbidden neighboring term is zero: its coefficient contains the
corresponding $[0]_v$. Comparing $(n,\beta)=(h+1,0)$ with $(1,h)$,
the diagonal coefficients are identical and the off-diagonal coefficients
differ by $(-1)^h$. Consequently the explicitly defined map
\begin{equation}\label{eq:certIraw}
 F_Ib_s=(-1)^{hs}b'_s
 \quad\text{satisfies}\quad F_IP_N=P'_2F_I,\qquad F_IQ_N=Q'_2F_I.
\end{equation}
Here the target has exterior labels $(l,m+h,c,d+h)$.
The definition is a cup-basis map, not an identity fitted to Racah entries.

The actual metric makes both radial operators selfadjoint. The following
elementary fact supplies the single metric correction.
\begin{lemma}[Bidiagonal metric uniqueness]\label{lem:gram}
Let $e_\ell,\ldots,e_u$ be a real basis and
\[
 Pe_k=\alpha_ke_k+\beta_ke_{k+1},\qquad
 Qe_k=\gamma_ke_k+\delta_ke_{k-1},
\]
with endpoint terms omitted. Suppose $\alpha_u\ne\alpha_k$ for $k<u$
and $\delta_k\ne0$ for $k>\ell$. A bilinear form making $P,Q$
selfadjoint is determined by its value on $(e_u,e_u)$.
\end{lemma}
\begin{proof}
Subtract two such forms and call the difference $H$, with
$H(e_u,e_u)=0$. Selfadjointness of $P$ gives
\[
 (\alpha_u-\alpha_k)H(e_u,e_k)=\beta_kH(e_u,e_{k+1}).
\]
Descending induction kills the top row. If row $k+1$ is zero, then
\[
 \delta_{k+1}H(e_k,w)
 =H(e_{k+1},Qw)-\gamma_{k+1}H(e_{k+1},w)=0
\]
for every $w$. Descending induction kills every row. Neither symmetry
nor positivity is needed for uniqueness.
\end{proof}

For the source in \eqref{eq:certIcoefficients}, the top index is
$u=\min(l,m)$ and the left diagonal difference is
\[
 A_{h+1,0}(u,m-u)-A_{h+1,0}(s,m-s)
 =v^{-l-m-h+1}[u-s]_v[l+m+h+1-u-s]_v>0
 \quad(s<u).
\]
The right lowering coefficient is nonzero for
$s>\max(0,m-c)$, since its two quantum-integer factors are $[s]_v$
and $[c-m+s]_v$. Thus the required top-eigenvalue separation and
nonzero interior lower coefficients hold, also when the index set has
only one element. Apply the lemma to the physical
metric and the pullback through $F_I$. The pullback equals $\kappa$ times
the source metric, where $\kappa=\|F_I b_u\|_2^2/\|b_u\|_N^2>0$.
Thus $E_I=\kappa^{-1/2}F_I$ is an onto isometry simultaneously
intertwining both radial operators. Their full eigenspaces are proved to
be exactly the actual irreducible coupling channels, and their channel
eigenvalues are distinct. Hence $E_I$ matches both complete channel bases
with the Type-I labels in Theorem~\ref{thm:physical}.
For $s,t\in\mathcal J(l,m,c)$ the actual channel labels and their shifts are
\begin{equation}\label{eq:certIlabels}
 i=l+m-2s,\qquad j=c-m+2t,\qquad i'=i+h,\quad j'=j+h.
\end{equation}
After the independent rank-two occupation-to-standard identification, the
map is the one called \lean{fullPhysicalHighestToStandard} in Lean.
This proves Type I on the whole physical space, not merely on a cup subspace.

\subsection{Type II: the full shifted-channel comparison}
The target is the standard rank-two highest space in
$S_{a+h}\otimes S_{b+h}\otimes S_c$ with output $S_d$.
Its left paths have contraction indices
\begin{equation}\label{eq:certIIpaths}
 (r+h,c-r),\quad 0\le r\le B,
 \qquad\text{and its right paths have indices}\quad(c-t,t+h).
\end{equation}
The triangle inequalities show that these lists exhaust the fixed-output
left and right paths. In particular their labels are
\begin{equation}\label{eq:certIIlabels}
 i=a+b-2r,\quad i'=i,\qquad
 j=b-c+2t,\quad j'=j+h.
\end{equation}
Let $L'_r,R'_t$ denote the independent standard unit CG trees for these
paths. Both are complete orthonormal bases. Define the full-space map
\begin{equation}\label{eq:certIIE}
 E_{a,b,c,d}L_r=(-1)^rL'_r.
\end{equation}
This is an onto isometry because it sends one independently constructed
orthonormal basis to another. This fact alone says nothing about the
right channels. The remaining steps establish that missing comparison.

\subsubsection*{The actual norm and the coordinate identity}
Let $V_r$ be the unnormalized physical bent left highest vector. The
physical cup, the actual partial-channel isometry and the pivotal trace give
\begin{equation}\label{eq:certIInorm}
 \|V_r\|_N^2=
 \frac{\dim_v T_{a+b,r}}{\dim_v S_d^{(N)}}.
\end{equation}
Here $T_{a+b,r}$ denotes the channel of Young shape $(a+b-r,r)$.
Both dimensions are actual pivotal traces, whose evaluation gives the
closed factorial expressions used below.
For the raw rank-two highest coupling of sizes $A,B$ and contraction $u$,
write
\[
 H(A,B,u)=v^{-2u(B-u)}
 \frac{\qf u\qf{A-u}\qf{B-u}\qf{A+B-u+1}}
      {\qf A\qf B\qf{A+B-2u+1}}.
\]
This is its independently evaluated squared norm in the standard compact
metric. Set $i=a+b-2r$ and
\[
 A_r=v^{-hr}\sqrt{\frac{H(a+h,b+h,r+h)}{H(a,b,r)}},
 \qquad B_r=\sqrt{\frac{[i+1]_v}{[d+1]_v}}.
\]
The explicit finite CG bending calculation, including its sign, proves that
the shifted raw bent vector equals $(-1)^rB_rL'_r$ after the fixed
weight-reversal identification. Direct cancellation of the norm factorials
then gives
\begin{align}
 A_rB_r&=\sqrt{C_{h;a,b,d}}\,\|V_r\|_N,\label{eq:certIIuniversal}\\
 C_{h;a,b,d}
 &=v^{-2hb}\frac{\qf h\qf a\qf b\qf{d+h+1}}
                      {\qf{a+h}\qf{b+h}\qf{d+1}}>0.
 \label{eq:certIIC}
\end{align}
This factor is independent of the intermediate channel $r$. It is the
factor for the occupation and bending normalizations used here; the
classical polynomial Fock map in \eqref{eq:scalarisometry} uses different
degree normalizations and has the factor $h!(h+1)!$.

Let $\pi_N$ retain physical occupation-basis terms supported on the first
two coordinates of each slot and identify each retained occupation with
its second-coordinate count. On the target, $\mathcal B$ reassociates the
tensor factors and reverses the third standard weight index $z\mapsto c-z$.
In these boundary coordinates $\xi_k^h$ extracts the coefficient
$(0,k+h,k)$ for $0\le k\le B$. In the unreversed standard basis that
coefficient has index $(0,k+h,c-k)$.
The complete finite CG calculation and \eqref{eq:certIIuniversal} imply
\begin{equation}\label{eq:certIIcoordinate}
 \xi_k^h(\mathcal B E_{a,b,c,d}x)
   =C_{h;a,b,d}^{-1/2}v^{-hk}\xi_k^0(\pi_Nx)
 \qquad\text{for every physical highest vector }x.
\end{equation}
To see the scaling on each $L_r$, write $W_r^h$ for the unnormalized
shifted bent vector in the boundary coordinates. The finite CG calculation
gives $A_r\xi_k^h(W_r^h)=v^{-hk}\xi_k^0(W_r^0)$.
Here $W_r^h=B_r\mathcal B E L_r$ and
$W_r^0=\|V_r\|_N\pi_N L_r$. Substitution of
\eqref{eq:certIIuniversal} gives \eqref{eq:certIIcoordinate} on each basis
vector, and linearity gives it on the whole space. The selected
coefficients of $W_r^h$ vanish for $k<r$ and have leading coefficient
\[
 \frac{(-1)^r}{\sqrt{H(d,c,r)H(a+h,b+h,r+h)}}\quad(k=r).
\]
Together with completeness this triangularity proves that the selected
coordinates determine every vector of the entire shifted highest space.
Thus equality of these coordinates will imply a full vector identity.

\subsubsection*{Creation transport, with the endpoints included}
Let $\mathcal C_N$ be the actual right cup creation map, increasing
$(b,c)$ to $(b+1,c+1)$, and let $\mathcal C'_2$ be the independent
shifted standard creation map. Write $f_k=\xi_k^0(\pi_Nx)$ and $g_k=\xi_k^h(\mathcal B Ex)$ in the
source degrees, with $\alpha=C_{h;a,b,d}^{-1/2}$. The two-term boundary
recurrence is intertwined by $g_k=\alpha v^{-hk}f_k$, since
\begin{align}
 &v^{-b-c+2k}g_k-v^{-h-1-b-c}g_{k-1}\notag\\
 &\hspace{16mm}=\alpha v^{-hk}
       \bigl(v^{-b-c+2k}f_k-v^{-1-b-c}f_{k-1}\bigr).
 \label{eq:certIIdifference}
\end{align}
At $k=0$ both previous-coordinate terms mean zero. Coordinates outside
the admissible finite ranges are also proved to vanish, so the same
argument covers an enlarged range after creation. Apply
\eqref{eq:certIIcoordinate} in the source and target degrees and use
the whole-space coordinate injectivity. This proves the actual operator
identity
\begin{align}
 E_{a,b+1,c+1,d}\mathcal C_N
   &=\gamma\,\mathcal C'_2E_{a,b,c,d},\label{eq:certIIcreation}\\
 \gamma&=\sqrt{\frac{C_{h;a,b,d}}{C_{h;a,b+1,d}}},
 &\gamma^2&=v^{2h}\frac{[b+h+1]_v}{[b+1]_v}.
 \label{eq:certIIgamma}
\end{align}
No right-channel identity was used to define $E$ or to prove this transport.

\subsubsection*{The Gram spectrum identifies every right channel}
For a mixed pair of degrees $b,c$, with contraction number $k$, the operator
$\mathcal C_N^*\mathcal C_N$ has eigenvalue
\begin{equation}\label{eq:certIIspectrum}
 \lambda_N(b,c;k)=v^{-2b-2c}
   \frac{[k+1]_v[b+c+N-k]_v}{[b+1]_v[c+1]_v},
 \qquad 0\le k\le\min(b,c).
\end{equation}
The eigenvalue is first computed on the explicit mixed highest vector.
Creation and its metric adjoint commute with the quantum-group action;
cyclic generation extends the equation to the whole actual irreducible
channel. The eigenvalues are distinct: successive numerator products have
difference $[b+c+N-2k-2]_v>0$ whenever both $k,k+1$ are admissible.
The denominators and the common power of $v$ are positive.
For the shifted rank-two mixed pair the label calculation and
\eqref{eq:certIIgamma} give exactly
\begin{equation}\label{eq:certIIspectralmatch}
 \gamma^2\lambda_2(b+h,c;k)=\lambda_N(b,c;k).
\end{equation}

Write primes for the shifted operators and use the adjoints given by the
source and target metrics. Since both degree maps $E$ are isometries,
\eqref{eq:certIIcreation} implies
\[
 E\,\mathcal C_N^*\mathcal C_N
   =\gamma^2(\mathcal C'_2)^*\mathcal C'_2 E,
\]
with $E$ in the original degrees. Every actual $R_t$ belongs to the mixed channel
$k=c-t$, and every independent $R'_t$ belongs to the shifted mixed channel
with the same $k$. These are separately proved membership statements.
Consequently $E R_t$ and $R'_t$ are unit eigenvectors of the same
transported Gram form with eigenvalue $\lambda_N(b,c;c-t)$.
Expand $E R_t$ in the complete orthonormal basis $\{R'_u\}$.
The eigenvector equation and the distinct eigenvalues kill every
coefficient except $u=t$; unit norm leaves a real sign. Hence
\begin{equation}\label{eq:certIIright}
 E R_t=\eta_tR'_t,\qquad\eta_t\in\{1,-1\},
 \qquad
 \langle L_r,R_t\rangle_N
    =(-1)^r\eta_t\langle L'_r,R'_t\rangle_2.
\end{equation}
One sign is attached to each whole column. This completes the Type-II
physical comparison and proves \eqref{eq:physical} by the direct route.

\subsection{The finite rank-two evaluation is also certified}
The target CG maps are built by solving the actual highest-vector
recurrence and applying the lowering operator, then dividing by the
positive highest-vector norm. Their tensor-tree overlaps define the
rank-two Racah matrix. The certificate proves its factorial evaluation:
the actual adjacent Casimir gives a finite Jacobi recurrence; the
terminating $q$-Racah candidate satisfies that recurrence; an independently
computed boundary coefficient and the nonzero recurrence coefficients
give uniqueness. The finite $q$-Saalschutz and Sears transformations used
to rewrite this answer as the supported Racah factorial sum are proved
within Lean as well.
Thus the rank-two Racah evaluation cited in the analytic proof is not an
unproved input to the formal certificate. Combining that evaluation with
the physical comparisons above, the actual pivotal dimension formulas,
and Lemma~\ref{lem:norm} proves the literal four-triangle expression
\eqref{eq:AMSproved}, with the full support \eqref{eq:AMSkernel}.

\section{Precise phase conventions}\label{sec:phases}
The following conventions fix the signs of the coefficients. All square
roots of positive quantities use the positive root. In the standard
rank-two weight basis $e_k^{(L)}$, the raising operator sends
$e_k^{(L)}$ to $[k]_v e_{k-1}^{(L)}$ for $k>0$ and kills $e_0^{(L)}$.
The diagonal Cartan eigenvalue is $v^{L-2k}$ and the squared norm is
$w_L(k;v)=v^{k(k+1-L)}\qf{k}\qf{L-k}/\qf{L}$.
The coproduct $E\otimes K+\id\otimes E$ gives the highest vector
$\sum_{k=0}^u c_k e_k^{(A)}\otimes e_{u-k}^{(B)}$ for the embedding
$S_{A+B-2u}\to S_A\otimes S_B$, with $0\le u\le\min(A,B)$ and
\begin{equation}\label{eq:CGphase}
 c_0=1,\qquad
 c_{k+1}=-\frac{c_k[u-k]_v}
 {[k+1]_v\,v^{B-2(u-k-1)}}.
\end{equation}
Normalize by its positive compact norm and generate descendants with the
fixed lowering action. The standard left and right trees are compositions
of these CG embeddings. This fixes their phases without using AMS.

For Type I, compose the cup-space isometry with the standard rank-two
identification and denote the composite again by $E_I$. Take
\[
 \widehat L_s=E_I^{-1}L'_s,\qquad
 \widehat R_t=E_I^{-1}R'_t.
\]
Their unit norms and actual higher-rank channel membership are proved by
the simultaneous radial comparison. For Type II use the full-space
isometry \eqref{eq:certIIE} and put
\begin{equation}\label{eq:certIIgauge}
 \widehat L_r=(-1)^r L_r=E^{-1}L'_r,
 \qquad\widehat R_t=E^{-1}R'_t=\eta_tR_t.
\end{equation}
The last equality follows from \eqref{eq:certIIright}; it is not a definition
of an arbitrary entrywise sign. Both vectors are actual unit vectors in
the required irreducible channels.

With $\rho=(r_1+r_2+r_3+r_4)/2$, define the normalized overlap and the
Wigner coefficient by
\begin{equation}\label{eq:certWignerphase}
 \widehat U_{ij}=\langle\widehat L_i,\widehat R_j\rangle,
 \qquad
 \Phi_{N,T}(\boldsymbol r;i,j)
 =\frac{(-1)^\rho\widehat U_{ij}}{\sqrt{d_i e_j}}.
\end{equation}
Here $\rho=l+c$ in Type I and $\rho=a+b$ in Type II.
With this convention the equality $\Phi_{N,T}=K_NI_{N,T}$ is exact.
Equivalently the standard rank-two convention is
\[
 U'_{ij}=(-1)^{\rho'}\sqrt{[i'+1]_v[j'+1]_v}\{6j\}'_v,
 \qquad\rho'=\rho+h,
\]
which explains the factor $(-1)^h$ in \eqref{eq:rank6j}.

For other normalized physical vertices, write
$\widetilde L_i=\alpha_i\widehat L_i$ and
$\widetilde R_j=\beta_j\widehat R_j$, with unit phases.
With an inner product conjugate-linear in its first entry,
\[
 \widetilde U_{ij}=\overline{\alpha_i}\beta_j\widehat U_{ij},
 \qquad\widetilde\Phi_{ij}=\overline{\alpha_i}\beta_j\Phi_{ij}.
\]
For real compact conventions these phases are signs. In particular the
original Type-II $L_r,R_t$ satisfy
$\Phi^{\rm raw}_{rt}=(-1)^r\eta_t\Phi_{rt}$, where $r,t$ denote
the contraction-path indices of \eqref{eq:certIIindices}.
The certificate proves the existence of these whole-channel signs at
each positive parameter; it does not select a parameter-independent sign
for an unspecified external vertex convention. Any analytic continuation
uses one continued whole-channel convention, not new signs chosen for
individual entries.

\noindent\textbf{Scope of the gauge prescription.}
The original phase comparisons fix the exterior quadruple before choosing
row and column signs. The theorem below strengthens them: for real
$v>0$, $v\ne1$, one table on actual ordered highest-weight triples
reproduces both calibrated AMS matrices simultaneously, at every
$N\ge3$. Its Type-I trees are exactly calibrated; its two Type-II
trees have the same exterior sign, which cancels from their overlap.
Rank two uses its previously certified common CG table. These results
specify an actual common gauge; they do not identify it with an unspecified
external choice of vertices.

\subsection{Global vertices: complete certification of one shared actual table}
\label{sec:globalvertices}
Throughout this subsection $v>0$, $v\ne1$. A \emph{global vertex}
is chosen once for an ordered incident triple $(A,B;C)$ and reused in
every exterior quadruple containing that triple. A separate choice for
each quadruple, even if constant along a whole row or column, does not
establish this property.

\begin{theorem}[Exact coherence of highest carriers]\label{thm:globalcarrier}
Let $X,Y,Z$ be actual compact highest realizations with the same highest
weight and normalized highest generators. Write $J_{XY}$ for the
canonical full-algebra isometric identification of their cyclic
channels. Then
\[
 J_{YZ}J_{XY}=J_{XZ},\qquad J_{XX}=\id.
\]
Writing $P_X$ for the orthogonal projection onto the channel of $X$,
and $\iota_Y$ for channel inclusion, the full pair transfer satisfies,
for any fixed realization $C$ of that highest weight,
\begin{equation}\label{eq:globalcarrierfactor}
 T_{XY}=\iota_YJ_{XY}P_X
       =(\iota_YJ_{CY})(J_{XC}P_X).
\end{equation}
Thus changing the intermediate carrier introduces no residual sign.
\end{theorem}
\begin{proof}
Both $J_{YZ}J_{XY}$ and $J_{XZ}$ intertwine the entire algebra and send
the normalized highest generator of $X$ to that of $Z$. An intertwiner
on a cyclic module is determined by its value on the cyclic generator,
so these maps agree on the whole channel. The identity case is the same
argument. Substitution in the definition of $T_{XY}$ proves
\eqref{eq:globalcarrierfactor}.
\end{proof}
This proof, including actual mixed, reverse mixed and two-row carriers,
is checked in \lean{GlobalHighestCarrierCoherence.lean} and
\lean{GlobalMixedVertices.lean}. It proves coherence of the carrier
identifications; compatibility of the calibrated AMS coupling trees
with a common vertex table is an additional assertion.

\begin{lemma}[Global vertex signs have trivial holonomy]\label{lem:globalphases}
Assign $g(A,B;C)\in\{1,-1\}$ once to every ordered fusion triple, and
let $s(\tau)$ be the product of these signs over all vertices of a
labelled binary fusion tree $\tau$. The phase factor for a change from
$\tau$ to $\sigma$ is $s(\tau)s(\sigma)$. Along every finite path it
depends only on the endpoint trees; its product on every closed path is
one. In particular, on a three-input recoupling its factor is
\begin{equation}\label{eq:globalfourvertices}
 g(A,B;E)g(E,C;D)g(B,C;F)g(A,F;D).
\end{equation}
\end{lemma}
\begin{proof}
In a product along a path every intermediate tree sign occurs twice
and cancels because $s(\tau)^2=1$. Expanding the two three-input tree
signs gives \eqref{eq:globalfourvertices}.
\end{proof}
The arbitrary-path assertion is
\lean{GlobalVertexSigns.pathSign_closed}. It concerns the phase factors
along, for example, pentagon paths. It does not assert a new full
pentagon matrix identity or by itself identify the AMS gauges with a
global vertex family.

\begin{theorem}[Both AMS families share global vertices at rank two]
\label{thm:globalranktwo}
For $N=2$, the single actual CG table fixed by \eqref{eq:CGphase} realizes
the calibrated overlaps of both Type I and Type II, simultaneously for
every admissible exterior quadruple and channel label.
\end{theorem}
\begin{proof}
For Type I specialize the complete physical rank comparison to rank two
and zero background. For Type II specialize its full-space isometry to
$h=0$; both calibrated vectors map to the standard trees. Those trees
are, by construction, compositions of the same CG embeddings, each of
which depends only on its two input and output highest weights. The
contraction number is uniquely determined by these three labels.
\end{proof}
The two physical overlap statements are
\lean{physicalTypeI_global_vertices_rank_two} and
\lean{physicalTypeII_global_vertices_rank_two} in
\lean{AMSGlobalRankTwo.lean}.

\paragraph{Actual Type-II left vertices in every rank.}
Put $S=a+b=c+d$, and let the internal two-row representation be
$T=(S-r,r)$. Choose its fixed highest carrier in
$S_{S-r}^{(N)}\otimes S_r^{(N)}$.
Equation~\eqref{eq:globalcarrierfactor} factors the actual left pair
map through this carrier. Bending the $S_c$ input is natural with
respect to postcomposition:
\[
 \operatorname{bend}(AB)
   =(A\otimes\id)\operatorname{bend}(B).
\]
Consequently the \emph{entire actual left tree} factors into a vertex
for $(S_a,S_b;T)$ and a vertex from $S_d$ to $T\otimes S_c^*$.
Its calibrated normalization can be placed entirely at the latter
vertex, since it is
\begin{equation}\label{eq:globalIIleftscale}
 (-1)^r\left(\frac{\dim_v T}{\dim_v S_d}\right)^{-1/2}.
\end{equation}
For $N\ge3$ the two-row highest weight determines $S$ and $r$, so
this factor depends only on the three incident irreps. At $N=2$ we use
Theorem~\ref{thm:globalranktwo}. The complete tree factorization and
\eqref{eq:globalIIleftscale}, including zero labels, are checked in
\lean{GlobalTypeIILeftVertices.lean}.

\paragraph{An exact boundary sign and neighboring-degree coherence.}
Write $n=N-1\ge1$, and use $b\ge0$ for a cup background in this
paragraph. Let $G_{t,k}^{(n)}$ be the literal cup monomial with outer
degrees $l,c$, obtained by applying $k$ right creations followed by $t$
left creations to the pure seed of degrees $(l-t,b,c-k)$, where $t\le l$
and $k\le c$.
At the boundary
\[
 B_{t,k}^{(n)}=
 (l e_0,(t+k)e_0+b e_n,c e_0)
\]
the full coefficient is
\begin{equation}\label{eq:globalcupboundary}
 G_{t,k}^{(n)}(B_{t,k}^{(n)})=
 (-1)^{nk}v^{\,t(1-l-k)+k(1-n-b-c)}.
\end{equation}
Indeed, the pure first symmetric occupation forces every left creation
to use coordinate zero; the pure third symmetric occupation forces the
same for every right creation. The right creation coefficient at that
coordinate has sign $(-1)^n$. Multiplying the one-path coefficients
gives \eqref{eq:globalcupboundary}.

The unnormalized cup rank transport sends
$G_{t,k}^{(n)}$ to $(-1)^{(n-1)t}G_{t,k}^{(1)}$, with target background
$b+n-1$. For $t+k=m$, the positive power of $v$ in
\eqref{eq:globalcupboundary} is unchanged and the remaining boundary
factor is the single sign $(-1)^{(n-1)m}$, independent of the cup
term. This proves a boundary identity on the \emph{whole} cup span.
The same calculation gives
\[
 R_n C_L=(-1)^{n-1}C'_L R_n,
 \qquad R_n C_R=C'_R R_n
\]
between neighboring exterior degrees. These are literal creation
operators, not signs fitted to overlaps. The proofs are in
\lean{GlobalCupCreationCoherence.lean},
\lean{GlobalCupFirstBoundary.lean}, and
\lean{GlobalCupBoundaryTransport.lean}. Passing to the unitary rank
transport only introduces the already proved positive normalization
ratios, so the phase information is unchanged.

\paragraph{The calibrated Type-II right phase is fixed locally.}
For $a+b=c+d$, let $\beta_{abc}$ extract the coefficient at
\[
 (a e_0,b e_0,c e_0)
 \quad\hbox{in }S_a\otimes S_b\otimes S_c^*.
\]
Here the third occupation is the first-coordinate vector in the
opposite symmetric module. The certificate now proves
\begin{equation}\label{eq:globalIIpositive}
 \beta_{abc}(\widehat R_t)>0.
\end{equation}
To see this, take selected coordinate zero in the complete Type-II
boundary transport. Its multiplier is positive. On the shifted
standard tree, the highest first input leaves only the leading outer
CG term, and the remaining inner coefficient is the positive
last-weight factorial quotient. Thus the entire boundary coefficient
is strictly positive. Since the physical right-channel highest line is one-dimensional,
\eqref{eq:globalIIpositive} and unit norm uniquely determine its real
phase. These statements are
\lean{physicalTypeIICalibratedRight_first_boundary_pos} and
\lean{physicalTypeIICalibratedRight_unique_positive}.

There is also an exact phase comparison between neighboring
realizations of the same mixed internal representation. With
$r=c-t$, actual mixed creation satisfies
\begin{equation}\label{eq:globalIIrightstep}
 C_R\widehat R_t(a,b,c;d)=
 \sqrt{\lambda_{n;b,c;r}}\,
 \widehat R_t(a,b+1,c+1;d),
\end{equation}
where the channel index $t$ keeps its numerical value, and
\[
 \lambda_{n;b,c;r}=
 v^{-2b-2c}\frac{[r+1]_v[b+c+n-r+1]_v}
 {[b+1]_v[c+1]_v}>0.
\]
Actual creation maps the entire mixed channel to its successor, its
Gram form gives this normalization, and its zero boundary coefficient
is a positive multiple of the previous one. The unique positive unit
vector is therefore the next calibrated vector. This is
\lean{physicalTypeIICalibratedRight_creation} in
\lean{GlobalTypeIIRightStep.lean}. Independently,
\lean{GlobalMixedCreationCoherence.lean} proves on the \emph{entire}
irreducible channel that creation equals canonical highest transport
times an explicit positive scalar. Thus \eqref{eq:globalIIrightstep}
contains no hidden sign depending on the exterior first or output leg.

\paragraph{Whole-carrier comparison of the two mixed orientations.}
Write $M_{p,s}=L(p\omega _1+s\omega _{N-1})$ and choose its fixed
carrier as the contraction-zero channel in $S_p\otimes S_s^*$.
Its normalized highest vector is the pure tensor $H_p\otimes H_s^*$,
where $H_s^*$ denotes the highest vector of the opposite symmetric module.
In the reverse ordering it is $H_s^*\otimes H_p$. Let $J$ be canonical highest transport and let
$\beta_+,\beta_-$ extract the all-first-coordinate coefficients.
The new whole-carrier identity is
\begin{equation}\label{eq:globalorientationboundary}
 \beta_-(Jx)=v^{ps}\beta_+(x)\qquad(x\in M_{p,s}).
\end{equation}
To prove it, let $\epsilon_p$ extract the highest coefficient of $S_p$.
Then $\epsilon_pF_i=0$ and
$\epsilon_pK_i^{-1}=v^{-p\delta_{i0}}\epsilon_p$.
Projection of a lowering word onto the highest symmetric factor
therefore contributes $v^{-p c_0}$ in the forward ordering and no
factor in the reverse ordering, where $c_0$ counts occurrences of $F_0$.
In the opposite symmetric module a lowering word follows one occupation
path, whose all-first coefficient vanishes unless $c_0=s$.
This proves the identity on every lowering word. Such words span the
entire actual cyclic carrier, so linearity proves
\eqref{eq:globalorientationboundary}, including weight spaces with
multiplicities. No one-dimensional-weight-space assumption is used.
The proof is included in the full workstation replay in
\lean{MixedOrientationWordBoundary.lean},
\lean{HighestFunctionalTransport.lean}, and
\lean{GlobalMixedOrientationBoundary.lean}.

\paragraph{Type-I creation phases and the shared right outer vertex.}
For either calibrated Type-I channel vector with index $t$, its
all-first boundary coefficient has sign $(-1)^{hm+m-t}$.
At the relevant standard boundary the CG coefficient is a one-term sum
of sign $(-1)^{m-t}$. The whole-cup boundary transport contributes
$(-1)^{hm}$. Channel membership, unit norm and the strict signed
boundary inequality uniquely determine each vector. Actual creation
then satisfies
\begin{align}
 C_L\widehat L_t(l,m,c)
   &=(-1)^h\sqrt{\lambda_{n;l,m;t}}\,
       \widehat L_{t+1}(l+1,m+1,c),\notag\\
 C_R\widehat R_t(l,m,c)
   &=\sqrt{\lambda_{n;m,c;m-t}}\,
       \widehat R_t(l,m+1,c+1).
 \label{eq:globalIsteps}
\end{align}
The proofs include zero exterior degrees and use actual channel images
and creation Gram forms. They are included in the full replay in
\lean{GlobalTypeILeftStep.lean} and
\lean{GlobalTypeIRightStep.lean}. Iteration removes the choice of
mixed-pair realization from the two outer highest vectors.

More precisely, for $a+p=s+d$ and $s\le a$ there is an actual
isometric intertwiner
\[
 F_{a;p,s;d}:S_d\longrightarrow S_a\otimes M_{p,s},
\]
defined by the contraction-zero calibrated Type-II right tree and
projection onto the fixed carrier. Let $C_r^+,C_r^-$ be the normalized
forward and reversed mixed embeddings into degrees $(p+r,s+r)$ and
$(s+r,p+r)$, respectively. The same $F$ satisfies
\begin{align}
 (\id\otimes C_r^+)F_{a;p,s;d}H_d
   &=\widehat R_s^{\rm II}(a,p+r,s+r;d),\notag\\
 (\id\otimes C_r^-)F_{a;p,s;d}H_d
   &=(-1)^{hs}\widehat R_s^{\rm I}(a,s+r,p+r;d).
 \label{eq:globalsharedright}
\end{align}
At $r=0$, \eqref{eq:globalorientationboundary} and the strict boundary
signs determine the relative phase. Positive creation transport extends
the equality to every $r$. The composed maps intertwine the entire
quantum algebra and preserve the actual compact metrics; the Type-II
factorization holds on the whole output module. These results are included in the full
workstation replay in
\lean{GlobalTypeIIRightVertices.lean},
\lean{GlobalTypeIRightIIComparison.lean}, and
\lean{GlobalSharedRightVertices.lean}.
Thus the two families literally reuse one right outer vertex, with a
relative sign depending only on its incident labels.

\begin{theorem}[One shared actual AMS vertex table]\label{thm:actualAMStable}
Let $N=h+2\ge3$, $v>0$, and $v\ne1$. There is one table of real
isometric intertwiners
\[
 \Gamma_{AB}^{C}:V_C\longrightarrow V_A\otimes V_B
\]
on the ordered actual highest-weight triples occurring in the two AMS
families. The carriers $V_A$ and the table are fixed before choosing
an exterior quadruple or a family. Evaluating the two composed trees on the normalized output
highest vector gives, after a common external isometric identification,
\begin{align}
 L_s^\Gamma&=\widehat L_s^{\rm I},&
 R_t^\Gamma&=\widehat R_t^{\rm I}
 &&\text{in Type I},\notag\\
 L_r^\Gamma&=(-1)^{hc}\widehat L_r^{\rm II},&
 R_t^\Gamma&=(-1)^{hc}\widehat R_t^{\rm II}
 &&\text{in Type II }(a+b=c+d).
 \label{eq:actualAMStrees}
\end{align}
Consequently its actual recoupling matrices equal the calibrated AMS
matrices in both families, for every admissible label. With the
exterior Wigner sign and actual pivotal dimensions prescribed above,
its $6j$ coefficients satisfy the literal signed four-triangle
AMS Eq.~(35).
\end{theorem}

\paragraph{A table indexed by actual weights.}
For each actual dominant Dynkin weight $\lambda$, choose once a complete
finite compact highest carrier $V_\lambda$ and a normalized highest
vector $H_\lambda$. Identify any other realization of that weight with
it by the highest-preserving isometry of
Theorem~\ref{thm:globalcarrier}. The Lean index
\lean{ActualIrrep} is the Dynkin weight together with a
\emph{propositional} existence witness; no construction label is part
of the index. Thus $M_{p,0}=S_p$, $M_{0,s}=S_s^*$, the zero-degree
objects are one unit object, and, writing
\[
 T_{S,r}=L((S-2r)\omega _1+r\omega _2),\qquad 2r\le S,
\]
we have $T_{S,0}=S_S$ and $T_{S,r}=M_{S-2r,r}$ at $N=3$.
The native two-row carrier in $S_a\otimes S_b$ has this same actual
index with $S=a+b$. These equalities are checked in
\lean{ActualAMSCarriers.lean}, \lean{ActualTwoRowDynkin.lean}, and
\lean{ActualAMSVertexAliases.lean}.

\paragraph{One common positive boundary.}
Put $\epsilon_A(x)=\langle H_A,x\rangle$. Fix a real linear functional
$\Lambda_B$ using only the actual weight of $B$, in the following
order: on a pure symmetric carrier use occupation augmentation
(the sum of all occupation-basis coefficients); on a nonzero pure
opposite symmetric carrier use the same augmentation in its native
opposite basis; on $M_{p,s}$ with $p,s>0$ use the all-first coefficient
in its contraction-zero carrier inside $S_p\otimes S_s^*$.
Transport these functionals canonically to the chosen actual carrier.
The unit uses the first case. On the projected Cartan lines used
below, the pure mixed endpoint boundary agrees with this choice.
No independent functional is assigned to a second name for the same
weight; equality of augmentation and a single coefficient on the
entire carrier is not required.
Only these second-input objects occur in the six AMS vertex roles.

For an admissible triple, define $P_{AB}^{C}$ by
\begin{equation}\label{eq:actualpositivevertex}
 \|P_{AB}^{C}H_C\|=1,\qquad
 (\epsilon_A\otimes\Lambda_B)(P_{AB}^{C}H_C)>0,
 \qquad P_{AB}^{C}\text{ intertwines }U_v(\mathfrak{sl}_N).
\end{equation}
The relevant highest line is one-dimensional. Its two unit vectors
have opposite nonzero boundary values, so precisely one is positive.
A full intertwiner from the cyclic source is determined by this
highest image. Hence \eqref{eq:actualpositivevertex} chooses a
\emph{unique full map}, independent of its construction. Existence
and full isometry are established below for every AMS role; they are
not hypotheses left to the reader. The common boundary and this
uniqueness argument are formalized in
\lean{ActualAMSBoundary.lean}, \lean{ActualPositiveVertexTable.lean},
and the endpoint comparisons in
\lean{PureMixedCarrierBoundary.lean} and
\lean{ActualMixedBoundaryCartan.lean}.

\paragraph{The intrinsic phase.}
For $\lambda=(\lambda_1,\ldots,\lambda_{N-1})$ set
\[
 b(\lambda)=\sum_{i=1}^{N-1}i\lambda_i,\qquad
 D(A,B;C)=b(A)+b(B)-b(C).
\]
The shared table is explicitly
\begin{equation}\label{eq:actualAMSphase}
 \Gamma_{AB}^{C}=\gamma(A,B;C)P_{AB}^{C},\qquad
 \gamma(A,B;C)=
 \begin{cases}
  1,& B=S_k\text{ for some }k\ge0,\\
  (-1)^{hD(A,B;C)/N},&\text{otherwise}.
 \end{cases}
\end{equation}
On every relevant triple $D$ is a nonnegative multiple of $N$.
Thus this is a sign depending solely on the ordered actual triple.
In particular, $b(M_{p,s})=p+(N-1)s$ and $b(T_{S,r})=S$.
The complete list is
\begin{center}\small
\begin{tabular}{@{}llll@{}}
\toprule
$A,B;C$ & condition & $D/N$ & $\gamma$\\
\midrule
$S_{p+r},S_{s+r}^*;M_{p,s}$ & $p,s,r\ge0$ & $r$ & $(-1)^{hr}$\\
$S_{s+r}^*,S_{p+r};M_{p,s}$ & $p,s,r\ge0$ & $r$ & $1$\\
$M_{p,s},S_c;S_{p+c-s}$ & $s\le c$ & $s$ & $1$\\
$S_a,M_{p,s};S_d$ & $a+p=s+d,\ s\le a$ & $s$ & $(-1)^{hs}$\\
$S_a,S_b;T_{a+b,r}$ & $r\le\min(a,b)$ & $0$ & $1$\\
$T_{c+d,r},S_c^*;S_d$ & $r\le\min(c,d)$ & $c$ & $(-1)^{hc}$\\
\bottomrule
\end{tabular}
\end{center}
For the last two rows occurring together, also
$r\le\min(a,b,c,d)$ and $a+b=c+d$.
The $B=S_k$ clause includes $k=0$ and takes precedence when
$S_0=S_0^*$. Formula~\eqref{eq:actualAMSphase}, rather than the choice
of a row in this display, defines the phase at every coincident triple.
Its role evaluations are proved in
\lean{AMSDynkinVertexPhase.lean},
\lean{AMSDynkinPhaseMixedRoles.lean}, and
\lean{ActualTwoRowDynkin.lean}.

\begin{proof}[Proof of Theorem~\ref{thm:actualAMStable}]
We first identify the six native positive maps with entries of
\eqref{eq:actualpositivevertex}. The forward and reverse mixed embeddings
$C_r^+$ and $C_r^-$ have unit highest images by normalized highest
transport. After pairing the first factor with its highest vector,
occupation augmentation on the second factor is strictly positive. For the symmetric pair inclusion into $T_{a+b,r}$,
projection onto the highest first factor leaves the positive leading
normalized two-row coefficient. Their highest spaces have dimension
one by the complete mixed and symmetric fusion decompositions.

The two mixed outer maps are obtained by projecting the
contraction-zero Type-I left tree and Type-II right tree onto
$M_{p,s}$. Call them $G:S_{p+c-s}\to M_{p,s}\otimes S_c$ and
$F:S_d\to S_a\otimes M_{p,s}$, respectively. Their highest norms are one. Projecting the first
factor onto its highest line and using the established strict boundary
signs proves positivity in \eqref{eq:actualpositivevertex}.
For $F$ the mixed boundary is the positive coefficient of the
contraction-zero Type-II right tree. On pure endpoint carriers, projection
onto the first highest factor fixes a Cartan weight of the second factor.
The projected vector lies on the single occupation line selected by
that weight, where augmentation agrees with the specified boundary
coefficient. The forward/reverse comparison
\eqref{eq:globalorientationboundary} identifies the relative Type-I
phase in \eqref{eq:globalsharedright}. Inserting the fixed mixed carrier into the actual
triple embeds each outer highest space injectively into its known
one-dimensional physical channel line. The nonzero unit highest
image proves that its dimension is exactly one.

For the remaining positive map
$O:S_d\to T_{c+d,r}\otimes S_c^*$, take the normalized bent two-row
pair projection with scalar \eqref{eq:globalIIleftscale}.
After pairing the two-row factor with its normalized highest vector,
the cup sum has exactly one nonzero occupation term,
$(c-r)e_0+r e_1$. Before that scalar the boundary is
\[
 (-1)^r\mathcal N_v(d-r,c-r,r)^{-1/2},
\]
where the already proved finite CG norm is
\[
 \mathcal N_v(u,w,r)=v^{-2rw}
 \frac{\qf{r}\qf{u}\qf{w}\qf{u+w+r+1}}
 {\qf{u+r}\qf{w+r}\qf{u+w+1}}>0.
\]
The squared norm $\mathcal N_v(u,w,r)=H(u+r,w+r,r)$ is
\lean{couplingNormClosed} in \lean{QuantumCouplingNormFormula.lean}. The boundary sign $(-1)^r$
cancels the sign in \eqref{eq:globalIIleftscale}, giving a positive
boundary. Applying $O$ and then including the two-row channel into
$S_a\otimes S_b$ gives the calibrated left highest vector. This
inclusion preserves the norm, so $O$ has a unit highest image. The same inclusion embeds its
highest space into the complete one-dimensional Type-II left channel.
This proves existence, uniqueness and positivity for the sixth role.
These are actual full intertwiners: canonical transport, projection,
tensor inclusion and bending all intertwine the entire presented
quantum algebra.

For completeness, isometry on the \emph{whole source} is checked
separately. Mixed embeddings are canonical highest isometries, and
the two-row map is isometric inclusion. Compose each outer map with
its isometric carrier inclusion into the corresponding physical
triple. The composite and the previously constructed full unit
highest-tree embedding intertwine the algebra and agree on the
highest generator. Cyclic uniqueness makes them equal on every
vector. The full tree embedding and the carrier inclusion preserve
the compact forms, so the outer map does too. This proves full
isometry for all six roles, not only normalization of a highest
vector. The concrete proofs are in
\lean{AMSOuterVerticesFullIsometry.lean} and
\lean{SharedActualAMSUnitary.lean}.

Transport each native positive map to the chosen actual carriers.
The common positive boundary, one-dimensional highest line and cyclic
uniqueness identify it with the \emph{same} $P_{AB}^{C}$, whatever
native role supplies the construction. This proves compatibility
also at $N=3$, at zero labels and whenever an inner vertex is an
outer vertex in a different quadruple. Multiplying by
\eqref{eq:actualAMSphase} gives the six full-map identities in
\lean{SharedActualAMSVertexTableSpec}; they are proved by
\lean{amsActualVertexTable_shared}.

It remains to calculate the two tree phases. In Type I the left
mixed embedding contributes $(-1)^{hr}$, exactly the left creation
phase in \eqref{eq:globalIsteps}, and the left outer vertex has phase
one. The reverse inner vertex has phase one; the right outer phase
$(-1)^{hs}$ cancels the right comparison sign in
\eqref{eq:globalsharedright}. Both trees are therefore exactly their
calibrated counterparts. In Type II the same-orientation inner
vertex has phase one and $O$ has phase $(-1)^{hc}$. For the right tree
write $c=s+r$. Its inner and outer phases multiply to
$(-1)^{hr}(-1)^{hs}=(-1)^{hc}$. This proves the four identities
\eqref{eq:actualAMStrees} in every native realization.

Finally, changing an internal carrier inserts $J^{-1}J=\id$
between the two vertices of either tree. The only surviving
transport is the same external triple isometry for both trees.
Its preservation of the compact pairing, and
$((-1)^{hc})^2=1$, prove equality of the actual shared-table overlaps
with the calibrated overlaps. The already certified actual pivotal
dimensions and finite factorial evaluation then give AMS Eq.~(35)
for both families. No factorial candidate defines a vertex or a tree.
The four actual-tree identities, exact overlap comparison, and literal
four-triangle evaluations are respectively
\lean{ActualAMSQuadrilateralComparison.lean},
\lean{ActualAMSRecoupling.lean}, and
\lean{ActualAMSCoefficientCertificate.lean}.
\end{proof}

The combined formal entry point is
\begin{center}
\lean{AMSPhysical.AMS_Eq35_shared_actual_vertex_certificate}.
\end{center}
It fixes the concrete table before universally quantifying every
admissible exterior and channel label. The existence theorem
\lean{exists_shared_actual_unitary_AMS_vertex_table} also states
explicitly that every admissible entry is an isometric full
intertwiner. Neither theorem assumes a common table, a unitary outer
vertex, or a phase factorization as an unproved hypothesis.

\paragraph{Scope of global coherence.}
The actual phase \eqref{eq:actualAMSphase} specializes
Lemma~\ref{lem:globalphases}: its phase-transition product along every
finite closed path of admissibly labelled trees is one, as proved by
\lean{amsDynkinVertexPhase_closed_path}. Together with
Theorem~\ref{thm:actualAMStable}, this certifies one shared actual AMS
vertex table for both families at $N\ge3$; rank two is covered by
Theorem~\ref{thm:globalranktwo}. The actual-table construction is proved
for real $v>0$, $v\ne1$. The coefficient deformation limits and
compatible complex continuations in the next section remain broader
\emph{coefficient} statements; a separate classical or arbitrary-complex
shared tensor category is not asserted. Likewise the closed-path
phase theorem is not a new full pentagon matrix certificate for all
irreducible representations.

\section{The classical limit and generic complex deformation}\label{sec:continuation}
\subsection{The physical limit at \texorpdfstring{$v=1$}{v=1}}
The regular quantum integer is the finite Laurent polynomial
\[
 [k]_v=\sum_{a=0}^{k-1}v^{k-1-2a}\quad(k\ge1),
 \qquad [0]_v=0.
\]
It is analytic near every positive parameter and equals $k$ at one.
All supported factorials in \eqref{eq:AMSkernel} consequently tend to
ordinary factorials. The squared triangle factors are positive at one,
so their positive square-root branches are analytic there. Since the sum
is finite, the proved positive-real physical identity gives
\begin{equation}\label{eq:certclassical}
 \lim_{\substack{v\to1\\v>0,\ v\ne1}}\Phi^{\rm phys}_{N,T}(v)
       =K_N(1)I_{N,T}(1).
\end{equation}
In Lean the classical coefficient is defined by the limit of the actual
physical coefficient family, using \lean{limUnder}; the right-hand side
is then proved equal to that limit. A separate theorem identifies this
limit with the value at one of the independent finite complex CG-tree
continuation described next. The certified $v=1$ assertion is this physical
deformation-limit statement. It is not a separate formalization of the
classical $SU(N)$ representation category; the classical analytic proof
above supplies the usual representation-theoretic context.

\subsection{An independent physical complex coefficient}
For each fixed admissible exterior quadruple construct $F_{ij}(z)$ by
finite CG trees and the independently evaluated higher-rank dimensions.
More explicitly, the continued rank-two overlap is
\begin{equation}\label{eq:certcomplexCG}
 U'_{ij}(z)=\sum_{p,q,r}
   w_A(p;z)w_B(q;z)w_C(r;z)
   L'_i(p,q,r;z)R'_j(p,q,r;z),
\end{equation}
where the sums run over the finite occupation bases, $w$ are their
continued compact metric weights, and each tree coefficient is a finite
CG expression. This is a complex bilinear continuation of the real
metric overlap; no complex conjugation of the parameter is introduced.
To match the literal Lean definitions, write
$\mathcal R(w)=\exp(\tfrac12\operatorname{Log}w)$ for the initial
normalization representative. It agrees with the positive square root
on $w>0$ and is holomorphic on the slit plane; continuation subsequently
uses the matrix patches specified below. Set
\[
 F_{ij}(z)=\frac{(-1)^\rho U'_{ij}(z)}{\mathcal R(d_i(z)e_j(z))}.
\]
Finite Laurent quantum factorials and square roots continued from their
positive values define these functions near the entire positive real
axis, including one. Their holomorphy there is proved from the explicit
CG expressions. Independent physical rank reduction and dimension
identification prove $F_{ij}(v)=\Phi^{\rm phys}_{ij}(v)$ for $v>0$,
$v\ne1$, before evaluating the AMS sum.

For the independent formula family, let $\Theta_{ij}(z)$ be the
product of the four \emph{squared} triangle factors in \eqref{eq:KN},
with parameter $z$. The literal formal expression is
\begin{equation}\label{eq:complexgroupedroot}
 G_{ij}(z)=[N-1]_z![N-2]_z!\,
                \mathcal R(\Theta_{ij}(z))\,I_{N,T}(z).
\end{equation}
Thus one square root is taken after multiplying the four squared factors,
as in \lean{complexFourTriangleRadicand} and
\lean{complexAMSExpression}. On positive parameters this equals $K_NI_{N,T}$.
For complex parameters it specifies the initial representative exactly;
it must not be replaced by a product of independently chosen pointwise
principal square roots of the four factors. For the whole finite matrix let
\begin{equation}\label{eq:certbranch}
 \mathcal U=\operatorname{Comp}_{2}\!\left(
  \bigcap_{i,j}\bigl(\operatorname{Hol}(F_{ij})
                         \cap\operatorname{Hol}(G_{ij})\bigr)\right),
\end{equation}
where $\operatorname{Hol}(f)$ is the open locus at which the specified
function is holomorphic, and $\operatorname{Comp}_{2}$ denotes the connected
component containing $z=2$. The certificate proves that $\mathcal U$ is
open, connected and contains the entire positive real axis.
The positive-real identity already proved gives $F_{ij}=G_{ij}$ on
$2+2^{-n}$, which accumulates at $2\in\mathcal U$.
The proved identity principle therefore gives
\begin{equation}\label{eq:certcomplexidentity}
 F_{ij}(z)=G_{ij}(z)\qquad(z\in\mathcal U,\ \text{all }i,j).
\end{equation}

\subsection{All compatible continuation branches}
A matrix patch consists of one connected open domain and one holomorphic
function for every matrix entry. A continuation chain starts with the physical matrix on a nonempty
subdomain of $\mathcal U$. Each step passes to a patch with nonempty
intersection with the previous domain, and the two matrices agree
throughout that intersection.
This definition concerns one base family only; it contains neither the
AMS candidate nor an assumed physical/formula comparison.
Replacing the initial physical family by the equal AMS family in
\eqref{eq:certcomplexidentity} preserves the same initial patch and every
subsequent step. Induction on the chain in both directions proves
\begin{equation}\label{eq:certcontinuations}
 \operatorname{Cont}_{\mathcal U}(F)=\operatorname{Cont}_{\mathcal U}(G).
\end{equation}
In particular, for every complex $z$ and every whole-matrix value $V$,
$V$ is a value of the continued physical family at $z$ if and only if it
is a value of the continued AMS expression there, using the same patch
chains. This is the precise generic-complex coefficient theorem.

This statement allows monodromy and uses compatible normalization branches.
It does not assert that every complex point is regular or reachable, or
that independently chosen pointwise principal square roots give the
continued channel convention. It also does not implement an independent
higher-rank tensor category at arbitrary complex deformation. The certified
complex scope is analytic continuation of the actual positive-real
physical coefficient family. Arbitrary nonsymmetric exterior modules and
root-of-unity fusion quotients are outside the theorem.

\section{Lean certificate and reproducible verification}\label{sec:audit}
The unified theorem is
\begin{center}
 \lean{AMSPhysical.AMS_Eq35_complete}.
\end{center}
Its conclusion is the conjunction of six explicit propositions: the
Type-I and Type-II physical identities for every positive $v\ne1$, the
two actual physical limits at one, and equality of the two pairs of full
compatible complex continuation families. All ranks $N=h+2\ge2$ and all
admissible exterior and channel labels are universally quantified.
The Lean parameter named \lean{q} is $v$ in this note, so AMS's
half-power parameter is its square.

\begin{center}
\small
\begin{tabular}{>{\raggedright\arraybackslash}p{0.31\linewidth}>{\raggedright\arraybackslash}p{0.61\linewidth}}
\toprule
Mathematical step & Lean source or theorem (namespace \lean{AMSPhysical})\\
\midrule
Complete Type-I space and isometry &
 \lean{QuantumCupHighestComplete.lean},
 \lean{QuantumFullHighestRankReduction.lean}\\[3pt]
Actual Type-I channel reduction &
 \lean{PhysicalTypeIRankReduction.lean},
 \lean{physicalTypeI_AMS_Eq35_four_triangles}\\[3pt]
Type-II actual norm and scalar &
 \lean{QuantumTypeIILeftNormClosed.lean},
 \lean{TypeIIShiftedUniversalNormalization.lean}\\[3pt]
Full Type-II creation transport &
 \lean{typeIIShiftedUnitEquiv_creation}\\[3pt]
Complete Type-II right-channel comparison &
 \lean{TypeIIStandardShiftedRightBasis.lean},
 \lean{typeIIShiftedUnitEquiv_right}\\[3pt]
Exact physical Type-II formula and gauge &
 \lean{PhysicalTypeIICalibratedGauge.lean},
 \lean{physicalTypeII_AMS_Eq35_calibrated}\\[3pt]
Actual classical limits &
 \lean{AMSClassicalPhysicalValue.lean},
 \lean{ComplexPhysicalClassicalLimit.lean}\\[3pt]
Whole-matrix complex branches &
 \lean{ComplexAMSMatrixBranch.lean},
 \lean{PhysicalAMSGenericComplex.lean}\\[3pt]
Unified coefficient theorem & \lean{AMS35Certificate.lean}\\[3pt]
Actual carriers, common positive boundary &
 \lean{ActualHighestCarrier.lean}, \lean{ActualAMSBoundary.lean}\\[3pt]
Six actual full isometric vertex roles &
 \lean{SharedActualAMSVertexTable.lean}, \lean{SharedActualAMSUnitary.lean}\\[3pt]
Actual shared-table trees and coefficients &
 \lean{ActualAMSRecoupling.lean}, \lean{ActualAMSCoefficientCertificate.lean}\\[3pt]
Combined shared-table theorem & \lean{AMSSharedVertexCertificate.lean}\\
\bottomrule
\end{tabular}
\end{center}

\noindent\textbf{Full workstation replay, 25 November 2026, 02:14 CET.}
All 985 source modules were rebuilt on
\lean{jmartinpc} with 24 parallel jobs and warnings treated as errors:
3,346 theorems and 1,136 supporting declarations.
The complete transitive dependency audit covered all 4,482
declaration roots and found exactly the standard Lean axioms
\lean{propext}, \lean{Classical.choice}, and \lean{Quot.sound}.
There are no \lean{sorry}/\lean{admit} placeholders, custom axioms,
unsafe declarations, assumed AMS identities, or unproved physical
comparison, common-table or unitary-vertex hypotheses in these statements.
The replay includes the concrete shared actual vertex table, all six
full-map roles, all four tree comparisons and both exact coefficient
evaluations in Theorem~\ref{thm:actualAMStable}.
The verifier automatically gives the combined audit a 64 MiB thread
stack; this resource setting adds no mathematical assumption.
Lean is pinned to 4.24.0, commit
\lean{797c613eb9b6d4ec95db23e3e00af9ac6657f24b}, with Mathlib v4.24.0.

The Lean archive \lean{AMS_Eq35_Lean_coefficients_and_shared_vertices.zip} has SHA-256
\begin{quote}\small
\nolinkurl{61ab62f2133e063fe7b8fdca722501621d18c5037e8e5716ef127afffab0fd04}.
\end{quote}
The archive contains every source, the pinned toolchain files, the
verifier and the successful audit evidence. In its \lean{AMSPhysical}
directory, replay with
\begin{quote}\small\ttfamily
python3 verify.py --jobs 24 --lean /path/to/lean \textbackslash\\
\hspace*{8mm}--mathlib /path/to/mathlib
\end{quote}
The report \lean{verification/report.json} records the host, toolchain,
all source hashes, required main theorem names, exact declaration counts
and axiom union. The complete logs are
\lean{verification/replay.log} and \lean{verification/Audit.log}.
The verifier requires the combined shared-table theorem, both
shared-table coefficient identities, and the isometry theorem for every
admissible vertex. The scope of the shared table is given in
Section~\ref{sec:globalvertices}.
\end{revision}


\section*{Acknowledgements}
The question addressed here arose during collaborative work with
Eduardo Serrano-Ens\'astiga, Jonathan Louvet, and Thierry Bastin on
\emph{Partial transpose as a Racah transformation and applications to
symmetric multiqudit systems}.

\begin{revision}
\section*{AI-use declaration}
This work is declared \emph{AI-derived}. John Martin posed the question
and directed the work. GPT-6.1 Sol derived the proof and wrote the
manuscript. GPT-6.2 Sol was used to review the
manuscript. The formal results were checked with Lean~4.24.0 and
Mathlib~v4.24.0, as described in Section~\ref{sec:audit}.
\end{revision}

\begin{thebibliography}{9}
\bibitem{AMS} V.~Alekseev, A.~Morozov, and A.~Sleptsov,
Multiplicity-free $U_q(\mathfrak{sl}_N)$ $6$-$j$ symbols: Relations,
asymptotics, symmetries, Nucl. Phys. B \textbf{960}, 115164 (2020),
\href{https://arxiv.org/abs/1912.13325}{arXiv:1912.13325},
\href{https://doi.org/10.1016/j.nuclphysb.2020.115164}{doi:10.1016/j.nuclphysb.2020.115164}.
\bibitem{NRZ} S.~Nawata, P.~Ramadevi, and Zodinmawia,
Multiplicity-free quantum $6j$-symbols for $U_q(\mathfrak{sl}_N)$,
Lett. Math. Phys. \textbf{103}, 1389--1398 (2013),
\href{https://arxiv.org/abs/1302.5143}{arXiv:1302.5143},
\href{https://doi.org/10.1007/s11005-013-0651-4}{doi:10.1007/s11005-013-0651-4}.
\bibitem{NT} S.~Neshveyev and L.~Tuset,
Notes on the Kazhdan--Lusztig theorem on equivalence of the Drinfeld
category and the category of $U_q\mathfrak g$-modules,
Algebr. Represent. Theory \textbf{14}, 897--948 (2011),
\href{https://arxiv.org/abs/0711.4302}{arXiv:0711.4302},
\href{https://doi.org/10.1007/s10468-010-9223-9}{doi:10.1007/s10468-010-9223-9}.
\bibitem{FH} W.~Fulton and J.~Harris, \emph{Representation Theory:
A First Course}, Graduate Texts in Mathematics 129, Springer (1991),
Sec.~6.1, Eq.~(6.8) (Pieri formula),
\href{https://doi.org/10.1007/978-1-4612-0979-9}{doi:10.1007/978-1-4612-0979-9}.
\bibitem{KR} A.~N.~Kirillov and N.~Yu.~Reshetikhin,
Representations of the algebra $U_q(\mathfrak{sl}(2))$, $q$-orthogonal
polynomials and invariants of links, in \emph{New Developments in the
Theory of Knots}, ed.~T.~Kohno, World Scientific (1990), pp.~202--256,
\href{https://doi.org/10.1142/9789812798329_0012}{\nolinkurl{doi:10.1142/9789812798329_0012}}.
\end{thebibliography}
\end{document}
